\section{Regulatory pathways: what each permits, restricts, and leaves uncertain}
\label{sec:regulatory}

\textbf{This is a preliminary regulatory analysis. It is not legal advice.} No statement in it
has been reviewed by a qualified generic-drug, CMC, sterile-manufacturing or 503B professional.
All sixteen gates carry \texttt{status: UNCERTAIN} with a null reviewer and a null review date.
UNCERTAIN is never PASS, so \texttt{regulatory.evaluate\_all} returns \texttt{NO\_CONCLUSION} for
every strategy in every run quoted in this memo. Nothing below asserts that any architecture is
permitted, compliant, clinically deployable, or ready to manufacture a drug. Gates are hard
constraints: none is relaxed anywhere to improve an economic result, and the engine carries no
penalty-cost substitute for one.

\subsection{Approved-generic contract manufacturing}

\textbf{Permits.} Supply of an approved presentation requires an approved application and a site
named in it. The decision map (\texttt{docs/regulatory/approved\_generic\_cmo\_map.md}) records
four positions: hold an ANDA or 505(b)(2) NDA; be added as a contract fill-finish site to a
holder's application by supplement; take transfer of an application; or supply components or bulk
only. Owners are identifiable from the frozen 2026-09-01 Drugs@FDA and Orange Book snapshots,
where the two modelled presentations sit behind 26 injectable applications from 23 sponsors
(norepinephrine) and 19 from 13 (sodium bicarbonate). Which of them would contract, license or
compete is an unresolved commercial question this package does not answer.

\textbf{Restricts.} Sites must be registered and drugs listed under FD\&C Act 510 and 21 CFR Part
207, and site additions are governed by 21 CFR 314.70. Batch disposition is non-delegable under
21 CFR 211.22(a), an extramural facility is an extension of the manufacturer under 21 CFR
200.10(b), and a quality agreement cannot delegate CGMP responsibility. Nominally separate sites
under one application therefore share one quality unit and share its hazard, which revision R007
had to correct in five configurations that had asserted an independence their own
regulatory-owner field denied.

\textbf{Leaves uncertain.} The reporting category for adding a third-party sterile fill-finish
site is expert interpretation until a reviewer assigns it. No framework grants automatic
statistical interchangeability across sites, so gate G14 blocks any fleet-level equivalence claim
while UNCERTAIN, removing the pooled-validation and shared-quality-unit savings that are the whole
economic case of the modular fleet family. The distributed-manufacturing registration instrument
is only a proposed rule (FR doc 2026-14073, 91 FR 42888), the base case is forbidden to depend on
it, and a word search of its text returns no occurrence of sterile, aseptic, 503B, outsourcing
facilit, or real-time release.

\textbf{What changes in the model.} G01 excludes a strategy outright if the relationship is
absent, so every non-503B result in this memo presumes an application relationship that does not
exist. G02, G07, G13 and G16 are lead-time gates the engine serves with tier-5 stand-ins rather
than with a category read from a regulatory assessment. G03 to G05 set the earliest commercial
start, modelled as a one-time cost plus a delay and not as a probability of failure. G06 caps
shelf life, which binds for several architectures in the current threshold grid. G16 blocks an
alternate API or container source until its lots are qualified under 21 CFR 211.84 and 211.94 and
a Type II drug master file letter of authorisation issues; it applies to S9, S18 and S19, and S9
is one of only two architectures meeting the service target on both products, so its source
independence rests on an UNCERTAIN gate.

\subsection{503B outsourcing}

\textbf{Permits, and for exactly how long.} A registered outsourcing facility may compound from a
bulk substance only if that substance is on the 503B bulks list or the drug is on FDA's shortage
list, and 503B(a)(5) bars copying an approved product except through the shortage exception. The
predicate must hold at compounding, at distribution, and at dispensing. \textbf{Eligibility is
therefore tied to the shortage list or the bulks list, and it evaporates when a shortage
resolves.} When the listing lapses, compounding stops, and inventory already compounded may not be
distributed or dispensed under the exception after that date as the statute is written; whether
transition relief exists is an unresolved question for a reviewer. The package models 503B as a
time-varying legal state, never a durable pathway and never a bridge to an approved application.

\textbf{The frozen record.} On the 2026-09-02 snapshots, sodium bicarbonate 8.4 percent carries 19
shortage rows of which 15 are current, first posted 2017-03-01, and is explicitly NOT INCLUDED on
the final bulks list (88 FR 20531), so the shortage listing is its only predicate. Norepinephrine
1 mg/mL carries zero shortage rows and appears on neither bulks table, so the pathway is
unavailable for it today. Furosemide 10 mg/mL, a reserve candidate, carries 33 rows of which 30
are current.

\textbf{What changes in the model.} The engine computes
\texttt{eligible\_today = on\_shortage\_list AND static\_gates\_pass AND state\_permission} and
re-checks the listing on each compounding, shipment and issue day. The only 503B cells ever
feasible in this study were produced by the forbidden world in which the product is listed at time
zero and the resolution rate is set to zero, that is, by assuming permanence. G12 removes regions
without state licensure from the service set, but no state-by-state map is frozen here, so that
restriction is unmodelled rather than modelled as absent. The beyond-use date is also not a scalar
in the source: 6 days for aseptic product released without a completed sterility test, 14 for
terminally sterilized product on a validated indicator-monitored cycle, 28 for either route with a
completed passing test. That instrument is revised draft guidance stamped nonbinding and not for
implementation, so it is not in-force policy, and the model holds \texttt{bud\_503b\_days} constant
where the source is a lookup on route by test completion by storage condition.

\begin{table}[htbp]
\centering
\caption{The sixteen gates, their treatment in the model, and their legal-status label. Source:
\texttt{config/regulatory\_gates.yaml}. Every status is UNCERTAIN with a null reviewer and a null
review date.}
\label{tab:gates}
\footnotesize
\setlength{\tabcolsep}{5pt}
\begin{tabular}{@{}p{6.0cm}llp{3.2cm}@{}}
\toprule
Gate & Pathway & Treatment & Legal basis type \\
\midrule
G01 application or CMO relationship & approved & exclude if fail & law or regulation \\
G02 site in application, registered & approved & lead time & law or regulation \\
G03 process and cleaning validation & approved & hard feasibility & final guidance \\
G04 aseptic or terminal validated & approved & hard feasibility & final guidance \\
G05 analytical methods transferred & approved & hard feasibility & law or regulation \\
G06 stability and container closure & approved & cap & law or regulation \\
G07 change category and submission & approved & lead time & law or regulation \\
G16 alternate source qualified & approved & lead time & law or regulation \\
G08 outsourcing facility registered & 503B & exclude if fail & law or regulation \\
G09 bulk eligible at all three dates & 503B & dynamic state & law or regulation \\
G10 not barred as essentially a copy & 503B & dynamic state & final guidance \\
G11 stability, BUD and sterility & 503B & cap & law or regulation \\
G12 state licensure and distribution & 503B & region restriction & law or regulation \\
G13 per-unit registration supportable & distributed & lead time & proposed rule \\
G14 cross-site comparability accepted & distributed & no fleet claim & unresolved question \\
G15 release model validated & release & fallback conventional & expert interpretation \\
\bottomrule
\end{tabular}
\end{table}

\subsection{What neither pathway resolves}

G15 governs any release-assurance component, its treatment is fallback to conventional release,
and the R3 scenario stays disabled unless it passes, so no strategy here takes credit for a
model-assisted release decision. That constraint costs almost nothing, which is itself a result:
on the illustrative release inputs, removing the entire chemical release queue removes two of the
sixteen base release days and G15 in full removes zero, because the one component it may reduce
sits underneath two longer ones while sterility incubation is the critical path.

Neither pathway gives the study's subject standing to change a sterilization route. Route
conversion is a prior approval supplement under 21 CFR 314.70(b)(2)(iii), with 314.97(a) applying
the same regime to ANDAs, and parametric release is a second prior approval supplement. Both are
filed by an application holder, and a 503B facility can hold neither. That matters because route
conversion with parametric release is the only intervention found across fourteen architecture
families that removes the binding release constraint, and it is unavailable to a party holding no
application.

\section{Limitations}
\label{sec:limitations}

\subsection{Inputs, and the absence of any human evidence}

Eighty-seven of ninety-one registered parameters carry evidence tier 5, illustrative, and zero are
missing. That covers every demand, cost, capacity, yield, uptime, lead-time, price, shelf-life and
disruption input; the only sourced numeric inputs driving the runs are four release-time
components. No absolute cost, break-even price, capital requirement, fill rate, shortage-days-avoided
figure or incremental cost per shortage-day avoided in this memo is a statement about
sterile-injectable manufacturing.

Zero interviews have been completed against a protocol target of 25 to 30, zero of three required
independent reviews, zero of sixteen gate reviews. Every counterparty named anywhere is a published
third party cited by source id: nobody is a partner, no conversation is a letter of intent, no price
has been quoted by a purchaser, no regulatory opinion has been obtained. All eleven protocol
revisions carry \texttt{approved\_by: pending founder sign-off}, and since the verdict is stated
against mean fill at or above 0.99 in at least 0.90 of runs, and the feasible set changes at 0.98,
an unsigned threshold is not a formality. Nine of twenty-two definition-of-finished tests pass, and
eleven of the thirteen incomplete ones cannot be closed by any amount of code. The modelling has run
ahead of the evidence.

Two input artifacts limit what any product comparison can mean. The two presentations differ mainly
through one derived tier-5 number, status-quo utilization 1.245 against 0.778, and the package's
audit records the dossiers as near-clones on several axes, so most apparent product discrimination
is an artifact of one assumption. The second product also failed the study's own beachhead
verification on 2026-07-16, so it is a control rather than a second market case.

\subsection{Three retracted headline claims}

This is kept rather than tidied away, because it bounds how much confidence any current number
deserves. Revisions R009 to R011 fixed six defects that carried the architecture ranking: free
opening inventory scaled to each design's own stock policy (NEW-1), unmatched search spaces between
the frozen comparators and the new architectures (MD-3, MD-12), an infeasible tie-break with no
Monte Carlo tolerance that had bought a 9.26 million USD per year cost increase for 0.00004 of fill
(MD-11), a material buffer that scaled with the lead time (MD-1), an inert allocation policy
(NEW-2), and readiness and take-or-pay instruments with no service channel (MD-23, MD-24). Three
earlier conclusions were then withdrawn. That owned distributed nodes meet the target nowhere on
either product: false, they meet it for norepinephrine. That no frozen comparator meets it: false,
dual sourcing meets it for norepinephrine at 0.960. That bright-stock postponement is the strongest
shape on both products: withdrawn, because once opening inventory is charged the design falls to
0.630 on sodium bicarbonate, having carried roughly 1.2 million free units, about five years of the
structural capacity gap, inside a five-year window. The last three headline claims this study made
did not survive its own next round of fixes.

\subsection{Open model defects and structural limits}

Twenty-nine defects are registered in
\texttt{docs/design\_space/bottleneck\_decomposition.md} section 7 and
\texttt{docs/design\_space/novel\_architectures.md} section 5. Twenty-one are recorded as fixed
under revisions R003 to R011, one (MD-1) with a residual the register states rather than argues
away. Eight remain open (Table~\ref{tab:defects}); MD-15 and NEW-3 could change the ranking rather
than only the numbers. One inconsistency is carried rather than smoothed: MD-17 is recorded as fixed
by R008 in \texttt{protocol/revisions.csv} and also appears among the deferred items in the later
status table.

\begin{table}[htbp]
\centering
\caption{The eight open model defects and what each could still move.}
\label{tab:defects}
\footnotesize
\setlength{\tabcolsep}{5pt}
\begin{tabular}{@{}p{1.6cm}p{5.0cm}p{6.4cm}@{}}
\toprule
Id & Defect & What it could still move \\
\midrule
MD-1 residual & Under a forced outage a shorter lead can still show marginally more stockouts, because a deferred order lands when the supplier resumes & Every raw-material lead-time condition in this study, in either direction \\
MD-5 & The fill metric forgives anything delivered inside a 7-day backorder window & Fill by up to 0.016 in an unstable direction, and so any verdict near the requirement \\
MD-8 & Node capacity is quadratic in the capacity factor, since batch size and batches per year both scale & Any future node-scale result; a node at scale 1.0 would be 1.82 times the central plant \\
MD-15 & One product per network, and equal regional shares & The largest unquantified distortion in the cost comparison: 94.5 percent of one distributed-node cell's cost is borne by one presentation at 0.203 utilization \\
MD-18 & The frozen release-time design variable is read only by the deterministic screen, never by the stochastic engine & The release-layer variant returns fill identical to the plain node design while costing 1.0 and 1.7 million USD per year more \\
MD-22 & The screen's common-impact factor is derived from topology for new architectures and hand-set for frozen comparators & No deterministic-screen common-impact comparison across the two arms is valid \\
MD-25 & Shared operating-layer cost is charged only outside release scenario R0, and every design-space strategy is R0 & A declared shared operating layer carries hazard with no cost line and no benefit \\
NEW-3 & No readiness-decay hazard, so exercising a line buys production volume and not reliability & The whole readiness case for the rotating and standby designs, including the architecture ranked first below \\
NEW-4 & Every declared common-cause group draws its own Poisson stream at one global rate, so group count is not neutral & A design declaring an honest correlation is penalised for declaring it, distorting comparison between differently documented topologies \\
\bottomrule
\end{tabular}
\end{table}

Three further engine limits carry no id: material stock has no shelf life, retest period or hold
limit, so upstream reserves escape the obsolescence that emptied every real one; per-site
commissioning, activation lead, exercise cadence and activation failure probability are
configuration constants no design variable can reach; and the common-cause group type cannot
express a key starting material tier, a regulatory-action group, or a shared utility. The engine
also runs one presentation through one network, so every portfolio, postponement and pooling
statement here is an accounting proxy, and the pooling threshold the study computed, on the order
of five comparable presentations per node, is untestable inside its own engine. Regions are
identical by default: R011 added the plumbing for differing criticality weights and minimum
guarantees and measured that, with identical regions, three of the four allocation policies return
the same fill to nine decimal places.

A later and deeper optimization exists at 60 final runs (\texttt{opt\_20260906T043357Z}), with the
design-space and ablation batteries re-run on its designs, but those artifacts have not been
re-evaluated at 100 runs and no narrative document interprets them. They disagree with the verified
set: there 29 of 40 cells return optimal rather than 14, and the status quo itself clears the target
on norepinephrine at 12.91 million USD per year, the cheapest cell in the run. The documentation lags
the artifacts, and the feasible boundary is not sharp at any resolution the package has run.

\subsection{Evidence limits outside the engine}

\textbf{The sterilization-route data is redacted, not merely unpublished.} That distinction decides
whether more searching would help, and it does not. Across 65 manufacturer labels behind 50
applications for the six longlist candidates, zero state a sterilization method, and for the two
modelled presentations the screen determined nothing at any confidence: 0 of 24 labels and 0 of 18
applications. FDA writes the route down and withholds the word. The published NDA 214313 review reads
``The proposed drug product is (b) (4) sterilized'', and the Product Quality Review template's
``Method of Sterilization:'' field is withheld behind a (b)(4) redaction in the published NDA 207987
review. A Freedom of Information Act request returns the same redaction, because it is the holder's
trade-secret claim and not an agency backlog.

\textbf{No price evidence, weak backcasts, a methods-proof release layer.} A landscape review over
313 sources found no published rate card, minimum order value, reservation fee or per-unit price for
US sterile fill-finish, 503B product, enterprise quality systems, shortage-data platforms or modular
hardware, so every unit price and break-even figure here is an assumption. Reimbursement, market
price, willingness to pay and production cost are kept as distinct concepts; CMS average sales price
and Part B spending appear only as labelled utilization proxies, and no reimbursement figure in this
memo is a manufacturing cost. The backcasts compare three ongoing shortages of one class against a
requirement of four episode classes, and no documented single-site failure episode has been obtained.
The release-assurance evidence rests on 615 of the 655 tablets in a 2002 public near-infrared
benchmark, with no real lots, no real drift, no real operators and not the target chemistry; it says
nothing about a sterile injectable, and abstention is never release.

\textbf{Known defects in the study's own claims, and authorship.} Fourteen claim-register rows are
recorded as unsupported, contradicted or retracted, one of them a citation error inside the
regulatory material: the repository's release-constraints memo attributes to FDA's 2004 aseptic
processing guidance a sentence that is not in it. The conclusion below does not rest on that
sentence; the citation is still wrong. Every artifact in this package, including the code, the
configurations, the regulatory decision maps and this prose, was produced with an AI coding agent,
and the disclosure record marks essentially all of it as not yet reviewed by the author. The
legal-basis text in the gate configuration is that agent's classification and is unverified, and a
planned adversarial verification pass over the audit findings did not complete.

\section{Conclusion}
\label{sec:conclusion}

Stated once more so no sentence below is mistaken for a claim about the world: every number here is
engine behaviour under evidence tier 5 illustrative inputs, all sixteen gates are UNCERTAIN and
unreviewed, and no interview or independent review has taken place.

\subsection{Where regional manufacturing does not beat the alternatives}

Three results survive every revision of the engine, including the ones that overturned its other
conclusions.

A network whose saleable capacity is below demand cannot be stocked out of the deficit. On the
capacity-short product a full year of safety stock with daily base-stock review leaves the status quo
at mean fill 0.869 with a tail probability of zero, and on the current reversal grid the feasible set
is empty in all nine high-demand cells of each product, so at twice base demand no strategy of the
twenty meets the target anywhere. That region argues for capacity, not for regional capacity.

Regional nodes lose the days before they exist, and their conditions are tight. On the
capacity-short product the distributed node design needs commissioning at or below 741 days, a
supplier disruption rate at or below 0.127 against an assumed base of 0.200, shelf life at or above
17.6 months and changeover at or below 4.5 days, so it sits outside its own feasible region on the
supplier axis at the base assumption. Its cost there is the highest in the battery, 32.0 million USD
per year against 19.4 million for the cheapest cell that meets the target on that product.

Adding a validated release layer to those nodes buys nothing. The release-layer variant returns mean
fill identical to the plain node design to four decimals on both products while costing 1.0 and 1.7
million USD per year more, for two reasons, a model defect (MD-18) and a regulatory constraint (G15
held UNCERTAIN), neither of which is an argument for the layer.

\subsection{Where it clears the bar, and the two architectures that clear both products}

Regional manufacturing is not eliminated. On the product that is not capacity-short, the distributed
node design meets the frozen target at a tail probability of 0.940, which is not borderline, at 20.0
million USD per year. What had previously made it fail was in part an unequal search space: it had
been denied the inventory freedom its competitors were given, and that is the clearest case in the
package of a conclusion that was an artifact of the harness rather than a statement about
distribution. The honest form is narrow. Regional capacity can meet the target where the incumbent
plant is not capacity-short, at about 9 percent above rotating contracted campaigns, 11 percent above
plain dual sourcing and 32 percent above the cheapest feasible cell; it cannot meet it at competitive
cost where the plant is capacity-short, which is the case shortages actually present.

Of forty cells re-evaluated at 100 paired runs, fourteen meet the frozen target and only two
architectures meet it on \emph{both} products: rotating prequalified campaigns across three
already-registered third-party fill lines (S17), and dual finished-dose sourcing with an explicit key
starting material tier (S9). Neither is a plant the sponsor builds.

\begin{table}[htbp]
\centering
\caption{The two architectures meeting the service target on both products. Results at 100 paired
runs from \texttt{results/design\_space/reeval\_post\_R009.json}; conditions at 20 paired runs per
evaluation from \texttt{ds\_post\_R009/thresholds.csv}. All inputs tier 5 illustrative. Base demand
900,000 units per year (norepinephrine) and 1,200,000 (sodium bicarbonate); base common-cause rate
0.1 per year; base 45 batches per site-year; base changeover 3 days.}
\label{tab:clearboth}
\footnotesize
\setlength{\tabcolsep}{4pt}
\begin{tabular}{@{}p{2.3cm}p{1.7cm}rrrp{5.4cm}@{}}
\toprule
Architecture & Product & Cost & Fill & P(meet) & Binding conditions inside the swept ranges \\
 & & M USD/yr & mean & & \\
\midrule
S17 rotating campaigns & norepine. & 18.4 & 0.9975 & 1.000 & demand at or below 926,600 per year; changeover at or below 8.45 days \\
S17 rotating campaigns & sodium bicarb. & 19.4 & 0.9960 & 0.900 & demand at or below 1,191,000, below the base; common cause at or below 0.182; batches at or above 44.4 \\
S9 dual source, KSM tier & norepine. & 23.3 & 0.9975 & 0.920 & demand at or below 642,200, well below the base; common cause at or below 0.068; batches at or above 51.9; the other nine axes infeasible at both ends \\
S9 dual source, KSM tier & sodium bicarb. & 24.6 & 0.9987 & 0.940 & demand at or below 1,303,000; shelf life at or above 12.4 months; API lead at or below 125 days; common cause at or below 0.182; supplier rate at or below 0.385 \\
\bottomrule
\end{tabular}
\end{table}

Two features of Table~\ref{tab:clearboth} matter more than the ranking. Three of the four demand
ceilings sit at or below the base demand assumption, so the conditions and the base verdicts disagree
at the resolution of the grid. And for S9 on the resolved product, nine of twelve swept axes return
infeasible at both ends of their recorded range while the same cell is feasible at base at 100 paired
runs, which the package does not resolve. Read the conditions as an indication of which axes bind,
not as numbers to plan against. Neither architecture is contractable today: no purchaser has been
asked and no price has been quoted.

One condition applies to every architecture at once, and it is the strongest single result in the
current reversal grid. At base demand, raising the multi-site common-cause rate from its base of 0.1
to the top of its recorded range of 0.3 events per group-year collapses the feasible set from 20
strategies to 1 on norepinephrine and from 15 to 0 on sodium bicarbonate, and the one survivor is the
rotating campaign network. Against that, fixed quality and labour cost per node changes the preferred
strategy in zero of the eighteen demand-by-common-cause cells across both products, and neither it
nor node capital crosses a feasibility boundary in any of their 40 threshold rows. Correlated failure
decides this map. Fixed quality cost does not.

\subsection{What would overturn this}

\textbf{A demand denominator.} True utilization and installed capacity for the exact presentation
produces 40 of the 202 crossings on the current threshold grid, more than any other swept input, and
removes an architecture outright if it lands above that architecture's demand ceiling. Every ceiling
in Table~\ref{tab:clearboth} sits within a few percent of the base or below it.

\textbf{A price and a committed volume from a purchaser with authority.} On the pre-correction
feasible set of thirteen cells, five moved from commercially possible to impossible as the assumed
price fell from each design's own break-even to 18.00 USD per unit, which is an assumption and not
evidence. An answer of ``nobody funds standing availability'' leaves every capacity-carrying design
unfunded by construction.

\textbf{The G07 reporting category in writing.} If adding a third-party sterile fill site returns as
a prior approval supplement with a preapproval inspection, the fastest leg of the contracted family
lengthens, and the contracted family is where both surviving architectures sit.

\textbf{Closing the open defects and re-running.} Until NEW-3 gives readiness a decay hazard,
exercising a line buys production volume rather than reliability, and the architecture ranked first
above is routine contracted supply wearing a standby label. Until MD-15 allows more than one product
per network, every pooling statement is an accounting proxy. The study has already shown twice that
closing defects of this size changes which architecture wins.

\subsection{The conditional answer, stated plainly}

On this evidence, regional manufacturing beats safety stock, dual sourcing and additional centralized
capacity only under a narrow conjunction: the presentation's incumbent capacity is not short of
demand, demand sits below roughly the base assumption, correlated failure stays near the low end of
its range, commissioning completes inside the measured window, and a buyer will pay a premium over
contracted alternatives for something other than service, because the service is matched and the cost
is not. Of the alternatives, additional safety stock alone never meets the target on the
capacity-short product anywhere in the searched space, and once every comparator searches the same
inventory space it is no longer distinguishable from the status quo, returning identical cost, fill
and tail probability in both current runs; dual sourcing meets the target on one product and not the
other; and rotating contracted campaigns on already-registered lines meet it on both, more cheaply
than nodes on either.

The defensible conclusion is therefore not an architecture. It is that the binding constraints in this
model are capacity relative to demand, correlated failure across nominally separate sites and
suppliers, the non-negotiable sterility hold, and a set of regulatory gates no design choice reaches,
and that none of those four is improved by distributing manufacture as such. Nothing here says any
architecture is FDA-compliant, clinically deployable, or ready to manufacture a drug. Whether the
conclusion also holds for sterile-injectable manufacturing in the world is not something this study is
entitled to say, and it will not be until the inputs are real, the gates are reviewed by qualified
people, and the defects in Section~\ref{sec:limitations} are closed.
